% GENERATED FILE. Edit canonical lessons, then run npm run generate:books.
% canonical-source-hash: fnv1a64:adba07919b711676
% canonical-lessons: TE-C254-rastram, TE-C254-prantam, TE-C254-dikku, TE-C254-turpu, TE-C254-padamara

\chapter{State, Area, Direction, East, West}
\label{ch:te-state-area-direction-east-west}
\hlchaptermodality{254}

\begin{chapteropening}
\textbf{By the end of this chapter:} \emph{I can say a state, an area, a direction, east and west.}

Complete the last lesson of chapter 254: I can say a state, an area, a direction, east and west.
\end{chapteropening}

\section[\texorpdfstring{\te{రాష్ట్రం} (rāṣṭraṁ)}{రాష్ట్రం (rāṣṭraṁ)}]{\te{రాష్ట్రం} (rāṣṭraṁ) — a state}

\label{lesson:TE-C254-rastram}



\begin{quote}
Before the new one: say the Telugu for in the past, then the Telugu for dawn.
\end{quote}

\subsection*{You'll want to know: \te{రాష్ట్రం}}
\textbf{\te{రాష్ట్రం}} — \emph{rāṣṭraṁ} — \textquotedblleft{}a state\textquotedblright{}.

\begin{grammarlens}[title={What you've built}]
One more word for this chapter.
\end{grammarlens}

\subsection*{Your turn}
\begin{itemize}
\raggedright
  \item \emph{Say it:} \emph{rāṣṭraṁ}
  \item \emph{Say it:} \emph{rāṣṭraṁ}, once more
  \item \emph{Say it:} say \emph{rāṣṭraṁ}
  \item \emph{Recall:} say the Telugu for in the past, then the Telugu for dawn, then say \emph{rāṣṭraṁ} again
\end{itemize}

\begin{tcolorbox}[breakable,colback=teal!4,colframe=teal!35!black,title={Before you move on}]
What does \te{రాష్ట్రం} mean? (\textquotedblleft{}A state\textquotedblright{}.) Say it once more.
\end{tcolorbox}
\section[\texorpdfstring{\te{ప్రాంతం} (prāntaṁ)}{ప్రాంతం (prāntaṁ)}]{\te{ప్రాంతం} (prāntaṁ) — an area, a region}

\label{lesson:TE-C254-prantam}



\begin{quote}
Before the new one: say the Telugu for dawn, then the Telugu for a state.
\end{quote}

\subsection*{You'll want to know: \te{ప్రాంతం}}
\textbf{\te{ప్రాంతం}} — \emph{prāntaṁ} — \textquotedblleft{}an area, a region\textquotedblright{}.

\begin{grammarlens}[title={What you've built}]
One more word for this chapter.
\end{grammarlens}

\subsection*{Your turn}
\begin{itemize}
\raggedright
  \item \emph{Say it:} \emph{prāntaṁ}
  \item \emph{Say it:} \emph{prāntaṁ}, once more
  \item \emph{Say it:} say \emph{prāntaṁ}
  \item \emph{Recall:} say the Telugu for dawn, then the Telugu for a state, then say \emph{prāntaṁ} again
\end{itemize}

\begin{tcolorbox}[breakable,colback=teal!4,colframe=teal!35!black,title={Before you move on}]
What does \te{ప్రాంతం} mean? (\textquotedblleft{}An area, a region\textquotedblright{}.) Say it once more.
\end{tcolorbox}
\section[\texorpdfstring{\te{దిక్కు} (dikku)}{దిక్కు (dikku)}]{\te{దిక్కు} (dikku) — a direction}

\label{lesson:TE-C254-dikku}



\begin{quote}
Before the new one: say the Telugu for a state, then the Telugu for an area.
\end{quote}

\subsection*{You'll want to know: \te{దిక్కు}}
\textbf{\te{దిక్కు}} — \emph{dikku} — \textquotedblleft{}a direction\textquotedblright{}.

\begin{grammarlens}[title={What you've built}]
One more word for this chapter.
\end{grammarlens}

\subsection*{Your turn}
\begin{itemize}
\raggedright
  \item \emph{Say it:} \emph{dikku}
  \item \emph{Say it:} \emph{dikku}, once more
  \item \emph{Say it:} say \emph{dikku}
  \item \emph{Recall:} say the Telugu for a state, then the Telugu for an area, then say \emph{dikku} again
\end{itemize}

\begin{tcolorbox}[breakable,colback=teal!4,colframe=teal!35!black,title={Before you move on}]
What does \te{దిక్కు} mean? (\textquotedblleft{}A direction\textquotedblright{}.) Say it once more.
\end{tcolorbox}
\section[\texorpdfstring{\te{తూర్పు} (tūrpu)}{తూర్పు (tūrpu)}]{\te{తూర్పు} (tūrpu) — east}

\label{lesson:TE-C254-turpu}



\begin{quote}
Before the new one: say the Telugu for an area, then the Telugu for a direction.
\end{quote}

\subsection*{You'll want to know: \te{తూర్పు}}
\textbf{\te{తూర్పు}} — \emph{tūrpu} — \textquotedblleft{}east\textquotedblright{}.

\begin{grammarlens}[title={What you've built}]
One more word for this chapter.
\end{grammarlens}

\subsection*{Your turn}
\begin{itemize}
\raggedright
  \item \emph{Say it:} \emph{tūrpu}
  \item \emph{Say it:} \emph{tūrpu}, once more
  \item \emph{Say it:} say \emph{tūrpu}
  \item \emph{Recall:} say the Telugu for an area, then the Telugu for a direction, then say \emph{tūrpu} again
\end{itemize}

\begin{tcolorbox}[breakable,colback=teal!4,colframe=teal!35!black,title={Before you move on}]
What does \te{తూర్పు} mean? (\textquotedblleft{}East\textquotedblright{}.) Say it once more.
\end{tcolorbox}
\section[\texorpdfstring{\te{పడమర} (paḍamara)}{పడమర (paḍamara)}]{\te{పడమర} (paḍamara) — west}

\label{lesson:TE-C254-padamara}



\begin{quote}
Before the new one: say the Telugu for a direction, then the Telugu for east.
\end{quote}

\subsection*{You'll want to know: \te{పడమర}}
\textbf{\te{పడమర}} — \emph{paḍamara} — \textquotedblleft{}west\textquotedblright{}.

\begin{grammarlens}[title={What you've built}]
One more word for this chapter.
\end{grammarlens}

\subsection*{Your turn}
\begin{itemize}
\raggedright
  \item \emph{Say it:} \emph{paḍamara}
  \item \emph{Say it:} \emph{paḍamara}, once more
  \item \emph{Say it:} say \emph{paḍamara}
  \item \emph{Recall:} say the Telugu for a direction, then the Telugu for east, then say \emph{paḍamara} again
\end{itemize}

\begin{tcolorbox}[breakable,colback=teal!4,colframe=teal!35!black,title={Before you move on}]
What does \te{పడమర} mean? (\textquotedblleft{}West\textquotedblright{}.) Say it once more.
\end{tcolorbox}
